%======================================================================
\section{A Minimality-Preserving Normal Form}\label{sec:normal-form}

\AP We now apply Theorem~\ref{thm:main} to logics containing $\BCI$. We use
the axiom schemes
\[
\begin{array}{l@{\qquad}l}
\intro*\axB=(b\to c)\to(a\to b)\to a\to c, & \intro*\axBp=(a\to b)\to(b\to c)\to a\to c,\\
\intro*\axC=(a\to b\to c)\to b\to a\to c, & \intro*\axS=(a\to b\to c)\to(a\to b)\to a\to c,\\
\intro*\axW=(a\to a\to b)\to a\to b, & \intro*\axK=a\to b\to a,
\end{array}
\]
and the \kl{logics} $\intro*\BCI=\lclos{\{\axB,\axC,\axI\}}$,
$\intro*\BCK=\lclos{\{\axB,\axC,\axK\}}$, $\intro*\IL=\lclos{\{\axK,\axS\}}$,
$\intro*\CL=\IL\lplus((a\to b)\to a)\to a$ and Gödel--Dummett
logic~\cite{Dum59} $\intro*\LC=\IL\lplus((a\to b)\to c)\to((b\to a)\to c)\to c$.
A formula belongs to $\CL$ if and only if it is true under every Boolean
valuation.
\plot{Check that the implicational axiom of $\LC$ agrees with \cite{NM21}.}

\begin{lemma}[Soundness]\label{lem:sound}
If $L\supseteq\BCI$, then $\alpha\to\St(\alpha)\in L$ for every
formula $\alpha$.
\end{lemma}
\begin{proof}
We argue in the implicational fragment of the sequent calculus for
$\FLe$, which is complete for $\BCI$ (see~\cite{GJKO07}). For a \kl{compound occurrence}
 $u$, let $\Gamma_u$ be the multiset of the \kl{defining premises}
$\dprem{v}$ such that $v=u$ or $v$ is \kl{strictly below} $u$; for a \kl{leaf} $u$, let
$\Gamma_u$ be empty. We show, by induction from the \kl{leaves} upward, that
\[
\Gamma_u,\subf{\alpha}{u}\vdash \rv{u}\quad\text{if $u$ is positive},\qquad
\Gamma_u,\rv{u}\vdash\subf{\alpha}{u}\quad\text{if $u$ is negative}.
\]
For \kl{leaves}, these sequents are axioms. Suppose that $u$ is \kl{positive} and
compound. Then $u0$ is \kl{negative} and $u1$ is \kl{positive}, so the induction
hypotheses and the rule $({\to}\mathrm{L})$ yield
$\Gamma_{u0},\Gamma_{u1},\subf{\alpha}{u},\rv{u0}\vdash \rv{u1}$. Applying
$({\to}\mathrm{R})$ and then $({\to}\mathrm{L})$ with
$\dprem{u}=(\rv{u0}\to \rv{u1})\to \ev{u}$, we obtain $\Gamma_u,\subf{\alpha}{u}\vdash \ev{u}$.
Negative occurrences are treated symmetrically. Every formula of
$\Gamma_u$ is used exactly once, so no structural rule other than exchange
is needed. Finally, taking $u=\rt$ gives
$\dprem{u_1},\dots,\dprem{u_m},\alpha\vdash \rv{\rt}$, and hence
$\vdash\alpha\to\St(\alpha)$ by exchange and $({\to}\mathrm{R})$.
\end{proof}

\begin{corollary}\label{cor:many-logics}
Let $L$ be a \kl{logic} containing $\BCI$, such as $\BCI$, $\BCK$, $\IL$,
$\LC$ or $\CL$. Then $L\lplus\alpha=L\lplus\St(\alpha)$ for every formula
$\alpha$. If, moreover, $\alpha$ is \kl{minimal} in $L$, then so is
$\St(\alpha)$.
\end{corollary}
\begin{proof}
By Lemma~\ref{lem:sound}, $\St(\alpha)\in L\lplus\alpha$. Conversely, the
substitution $\decsub{\alpha}$ of Section~\ref{sec:translation} maps every \kl{defining premise} $\dprem{u}$ to
$\subf{\alpha}{u}\to\subf{\alpha}{u}$ and $\rv{\rt}$ to $\alpha$, so
$\alpha\in L\lplus\St(\alpha)$ by (C1)--(C3). If $\alpha$ is \kl{minimal}
in $L$, then $\alpha\in L$, hence $\St(\alpha)\in L$ by
Lemma~\ref{lem:sound}, and Theorem~\ref{thm:main} applies.
\end{proof}

\begin{corollary}[Order-$4$ normal form]\label{cor:normal-form}
Let $L_0$ be a \kl{logic} containing $\BCI$. If $L=L_0\lplus\Gamma$, where
every $\alpha\in\Gamma$ is \kl{minimal} in $\CL$, then
$L=L_0\lplus\{\St(\alpha)\mid\alpha\in\Gamma\}$, and every $\St(\alpha)$ is
\kl{minimal} in $\CL$ and of \kl{order} at most $4$.
\end{corollary}
\begin{proof}
The first two claims follow from Corollary~\ref{cor:many-logics} applied to
$L_0$ and to $\CL$. Since every \kl{body} has \kl{order} $2$, every
\kl{defining premise} has \kl{order} at most $3$, and hence
$\Ord(\St(\alpha))\le4$.
\end{proof}

The bound $4$ cannot be improved: by the next proposition, a classical
tautology of \kl{order} at most $3$ never axiomatizes a proper extension of
$\IL$, whereas $\LC$ is axiomatized over $\IL$ by a formula \kl{minimal} in
$\CL$ (Corollary~\ref{cor:LC}).

\begin{proposition}\label{prop:order3}
Every formula of $\CL$ of \kl{order} at most $3$ belongs to $\IL$.
\end{proposition}
\begin{proof}
A formula of \kl{order} at most $2$ has the form $p_1\to\dots\to p_k\to q$ with
variables $p_1,\dots,p_k,q$, and we read it as the rule ``from
$p_1,\dots,p_k$ infer $q$''. A formula of \kl{order} at most $3$ has the form
$F=H_1\to\dots\to H_n\to r$, where the $H_i$ are of \kl{order} at most $2$ and
$r$ is a variable. Let $V$ be the closure of the empty set under the rules
$H_1,\dots,H_n$. The valuation that makes exactly the variables in $V$ true
satisfies every $H_i$, so $r\in V$ because $F\in\CL$. Following the
construction of $V$ step by step, we can derive $r$ from $H_1,\dots,H_n$ by
modus ponens in $\IL$. Hence $F\in\IL$.
\end{proof}
